\documentclass[12pt]{article}
\usepackage[T1]{fontenc}
\usepackage[utf8]{inputenc}
\usepackage{lmodern}
\usepackage{textcomp}
\usepackage{enumitem}
\usepackage{geometry}
\geometry{margin=1in}
\begin{document}
\begin{center}
Answer Key - Constitutional Law - Undergraduate Level Assessment 4 \
\small (Answers are ordered exactly as in the question paper.)
\end{center}

\section*{Multiple Choice}
\begin{enumerate}[label=\arabic*.]
\item \textbf{Answer:} C
\\ \textit{Options:} A) It is likely to generate more individual rights and greater liberty.; B) It renders a community more genuine.; C) It opens the door to authoritarianism.; D) It improves the moral justification for the exercise of political power.
\\ \textit{Note:} Dworkin's notion of law as integrity emphasizes coherence and moral reasoning in legal interpretation, suggesting that it opposes authoritarianism by advocating for a principled approach to law that respects individual rights and democratic values.
\item \textbf{Answer:} B
\\ \textit{Options:} A) ordinances; B) federal statutes; C) executive orders; D) charters
\\ \textit{Note:} The correct answer is "federal statutes" because the Commerce Clause grants Congress the authority to create laws that regulate commerce between states and with foreign nations, thereby establishing the legal framework for federal statutes that govern these economic activities.
\item \textbf{Answer:} C
\\ \textit{Options:} A) The purpose of immunity is to protect foreign Heads of State from embarrassment; B) Immunity protects a State from being invaded by another; C) Immunity shields States from being sued in the courts of other States; D) The purpose of immunity is to offer impunity in respect of all crimes
\\ \textit{Note:} Sovereign immunity serves to protect States from legal actions in the courts of other States, ensuring that a State's sovereignty and independence are respected and not undermined by the judicial decisions of foreign jurisdictions.
\item \textbf{Answer:} A
\\ \textit{Options:} A) Because it attempts to give effect to personal choice over individual luck.; B) Because liberty is more important than equality.; C) Because a market economy is just.; D) Because the state is the best arbiter of equality between individuals.
\\ \textit{Note:} Dworkin supports liberal egalitarianism because it emphasizes the importance of personal choice in determining individuals' life outcomes, rather than leaving those outcomes to the whims of chance or luck, thereby promoting fairness and equality in society.
\item \textbf{Answer:} B
\\ \textit{Options:} A) Recognition of governments is very prevalent in contemporary practice; B) Recognition of governments has largely been replaced by functional Recognition; C) Government recognition is common in respect of rebel entities; D) Only democratic governments are recognised in contemporary practice
\\ \textit{Note:} The correct answer highlights that in contemporary international practice, states often prioritize functional recognition-acknowledging the ability of a government to effectively govern and fulfill international obligations-over formal recognition, which can be influenced by political considerations and legitimacy issues.
\end{enumerate}

\section*{True / False}
\begin{enumerate}[label=\arabic*.]
\setcounter{enumi}{5}
\item \textbf{Answer:} True
\\ \textit{Options:} A) True; B) False
\\ \textit{Note:} The statement is true because Robert Nozick's formula r x H encapsulates the principle that the severity of punishment should be proportional to the degree of responsibility for the act and the harm it caused, reflecting a foundational concept in legal philosophy regarding justice and accountability.
\item \textbf{Answer:} True
\\ \textit{Options:} A) True; B) False
\\ \textit{Note:} Hume's critique of natural law asserts that moral judgments are rooted in human emotions and sentiments rather than objective truths, thus implying that moral truths cannot be universally known or validated.
\item \textbf{Answer:} False
\\ \textit{Options:} A) True; B) False
\\ \textit{Note:} The statement is false because international law, particularly under the United Nations Charter, stipulates that the right to self-defense and the use of force in response to an armed attack is limited to situations that involve a significant level of force, rather than allowing for invasion based on any level of aggression.
\item \textbf{Answer:} True
\\ \textit{Options:} A) True; B) False
\\ \textit{Note:} The Badinter Commission, established to address the legal issues arising from the breakup of Yugoslavia, emphasized that the recognition of the newly independent republics was contingent upon their adherence to human rights and democratic principles.
\item \textbf{Answer:} True
\\ \textit{Options:} A) True; B) False
\\ \textit{Note:} The UN Human Rights Committee is indeed a treaty-based mechanism because it oversees the compliance of states with the International Covenant on Civil and Political Rights, ensuring the protection and promotion of human rights as outlined in that treaty.
\end{enumerate}

\section*{Long Form Response}
\begin{enumerate}[label=\arabic*.]
\setcounter{enumi}{10}
\item \textbf{Answer:} Durkheim's perspective on punishment emphasizes its role in reinforcing societal norms and values, suggesting that punishment serves not only to deter crime but also to express collective sentiments, including vengeance. Within the framework of constitutional law, this notion of vengeance can manifest in the balance between retributive justice and rehabilitation. Modern legal systems often grapple with this duality, as punitive measures may reflect societal calls for vengeance while also striving to uphold principles of fairness and human rights. The implications of this tension are significant; legal systems must navigate the desire for retribution with the need for a more restorative approach to justice that promotes social cohesion and prevents recidivism. Thus, understanding Durkheim's view on punishment can help inform contemporary debates on the efficacy and morality of legal practices.
\\ \textit{Note:} correct because it illustrates how Durkheim's view of punishment as a societal expression of collective sentiments, including vengeance, highlights the tension between retributive justice and rehabilitation in modern legal systems, thereby connecting these concepts to the principles of constitutional law.
\item \textbf{Answer:} John Finnis critiques Hume's conception of practical reason by arguing that it limits the role of human reason in determining what constitutes a worthwhile life. Hume posits that reason is merely a tool for achieving ends dictated by passion, thus undermining its capacity to define moral values or goals. In contrast, Finnis asserts that human reason can and should actively engage in discerning the fundamental goods that contribute to a fulfilling life, such as knowledge, friendship, and aesthetic appreciation. By emphasizing the rational pursuit of these intrinsic goods, Finnis positions practical reason as a guiding force in ethical decision-making, ultimately asserting that a worthwhile life is one that is informed by rational reflection and the pursuit of common human goods. This perspective not only critiques Hume's limitations but also reaffirms the importance of reason in ethical considerations within constitutional law.
\\ \textit{Note:} correct because it accurately captures Finnis's critique that Hume's view reduces reason to a mere instrument of passion, whereas Finnis argues for a more active role of reason in identifying and defining the fundamental goods essential for a worthwhile life.
\end{enumerate}

\vfill
\noindent\textit{End of Answer Key}
\end{document}